\section{Integration of Established Theories}
\label{sec:integration}
%%%%%%%%%%%%%%%%%%%%%%%%%%%%%%%%%%%%%%%%%%%%%%%%%%%%%%%%%%%%%%%%%%%%%%%%%%%%%%%

A unified framework must demonstrate that established theories emerge as special cases.
This section shows how classical, behavioral, institutional, organizational, and financial 
theories correspond to specific restrictions on the complementarity matrix $C$, 
the context variable $\Psi$, and their dynamics.
Formal derivations are provided in Appendix~\ref{app:limitingcases}.

\subsection{The Unification Logic}

The general framework posits three equations:
\begin{align}
  U_i &= U_i(x_i, x_{-i}, \Psi) && \text{(Utility depends on others and context)} \\
  C_{ij} &= \frac{\partial^2 U_i}{\partial x_i \partial x_j} && \text{(Complementarity structure)} \\
  \Psi_{t+1} &= f(\Psi_t, a_t) && \text{(Context evolves with behavior)}
\end{align}

Each established theory imposes \emph{restrictions} on this structure.
Relaxing different restrictions yields different theories.
The full framework relaxes all restrictions simultaneously.

\subsection{Classical Economics: The Separability Restriction}

\paragraph{Arrow--Debreu General Equilibrium.}
The foundational model of competitive markets \citep{arrowdeb} assumes:
\begin{enumerate}
  \item \textbf{Separable utilities:} $C_{ij} = 0$ for all $i \neq j$
  \item \textbf{Exogenous context:} $\Psi_{t+1} = \Psi_t = \bar{\Psi}$
  \item \textbf{No identity:} $U^{IDN} = 0$
\end{enumerate}

Under these restrictions, each agent's optimization is independent.
The coherence index $K = 1$ trivially---there are no cross-agent complementarities to align.
The First Welfare Theorem holds: equilibrium is Pareto efficient.

\paragraph{What Classical Economics Cannot Explain.}
The separability assumption ($C_{ij} = 0$) rules out:
\begin{itemize}
  \item Social preferences (fairness, reciprocity, altruism \citep{fehrfischbacher2003})
  \item Strategic complementarity and coordination games
  \item Multiple equilibria from network effects
  \item Institutional change and path dependence
\end{itemize}

The framework shows \emph{why} these phenomena require extensions:
they all involve $C_{ij} \neq 0$.

\subsection{Behavioral Economics: The Social Preference Extension}

\paragraph{Fehr--Schmidt Inequity Aversion.}
\citet{fehrschmidt} introduce utility interdependence:
\begin{equation}
  U_i = x_i - \alpha_i \sum_{j} \max(x_j - x_i, 0) - \beta_i \sum_{j} \max(x_i - x_j, 0)
\end{equation}

This generates non-zero complementarities:
\begin{equation}
  C_{ij} = \frac{\partial^2 U_i}{\partial x_i \partial x_j} \neq 0
\end{equation}

The key insight: $(\alpha_i, \beta_i)$ are not universal constants but depend on context.
Laboratory experiments show they vary with framing, group identity, and institutional design.
Thus: $\alpha_i = \alpha_i(\Psi)$, $\beta_i = \beta_i(\Psi)$.

\paragraph{Bolton--Ockenfels ERC.}
An alternative specification where utility depends on \emph{share} of total payoff:
\begin{equation}
  U_i = U_i\left(x_i, \frac{x_i}{\sum_j x_j}\right)
\end{equation}

This also generates $C_{ij} \neq 0$, but the structure is state-dependent.
Both models are special cases of the $(C, \Psi)$ framework with interdependent utilities.

\paragraph{What Behavioral Economics Adds.}
The move from $C_{ij} = 0$ to $C_{ij} \neq 0$ explains:
\begin{itemize}
  \item Ultimatum game rejections (violate Arrow--Debreu predictions)
  \item Cooperation in public goods games
  \item Gift exchange in labor markets
  \item Context-dependent fairness norms
\end{itemize}

See Appendix~\ref{app:limitingcases} for formal derivations.

\subsection{Identity Economics: The Temporal Extension}

\paragraph{Akerlof--Kranton Identity.}
\citet{akerlofkranton} introduce identity as a utility dimension:
\begin{equation}
  U_j = U_j(a_j, a_{-j}, c_j)
\end{equation}
where $c_j$ is the agent's social category.

This adds two types of complementarity:

\textbf{Temporal complementarity:}
\begin{equation}
  \frac{\partial^2 U}{\partial a_t \partial a_{t+1}} > 0
\end{equation}
Behaving consistently with identity today reinforces that identity tomorrow.

\textbf{Social complementarity:}
\begin{equation}
  C_{ij} > 0 \text{ for } c_i = c_j
\end{equation}
Members of the same category reinforce each other.

\paragraph{Bénabou--Tirole Moral Identity.}
\citet{benaboutirole} model self-signaling:
\begin{equation}
  U = v \cdot a - c(a) + \mu \cdot \mathbb{E}[\theta | a]
\end{equation}

Agents care about what actions reveal about their character $\theta$.
This creates strong intertemporal complementarity and explains taboos---
actions that destroy self-image regardless of material payoff.

\paragraph{Integration.}
In the FEPSDE framework, identity maps to $U^{IDN}$:
the utility from being who one believes oneself to be.
The 144-component structure explicitly includes identity utility
across all six dimensions and four time horizons.

\subsection{Growth Theory: The Endogenous Context Extension}

\paragraph{Aghion--Howitt Endogenous Growth.}
\citet{aghionhowitt} model growth through innovation:
\begin{equation}
  A_{t+1} = A_t \cdot (1 + g(n_t))
\end{equation}

Technology $A_t$ is part of context $\Psi$.
Innovation changes context \emph{endogenously}:
\begin{equation}
  \Psi_{t+1} = f(\Psi_t, n_t) \neq \Psi_t
\end{equation}

This violates the Arrow--Debreu assumption of fixed context.

\paragraph{Creative Destruction as Coherence Dynamics.}
Innovation disrupts existing complementarities:
\begin{enumerate}
  \item Innovation occurs: $\Delta \Psi > 0$
  \item Reference structure shifts: $C^*(\Psi_{old}) \to C^*(\Psi_{new})$
  \item Coherence falls: $K_t \downarrow$ (old patterns no longer fit)
  \item Adaptation: Agents learn new complementarities
  \item Re-coherence: $K_t \uparrow$ at new equilibrium
\end{enumerate}

\paragraph{Growth as Repeated Cycles.}
\begin{itemize}
  \item \textbf{Stagnation:} $K \approx 1$ (stable), but $Q$ doesn't grow
  \item \textbf{Growth:} $K$ fluctuates (creative destruction), but $Q$ rises over time
\end{itemize}

The framework reinterprets Schumpeterian dynamics as coherence dynamics.

\subsection{Institutional Economics: The Multiple Equilibria Extension}

\paragraph{Acemoglu--Robinson Institutions.}
\citet{acemoglu} argue that institutions determine economic outcomes.
In the framework:
\begin{equation}
  \Psi = (\text{Property rights, Rule of law, Political openness}, \ldots)
\end{equation}

Institutions determine the complementarity structure:
\begin{itemize}
  \item \textbf{Inclusive institutions:} $C_{ij} > 0$ (cooperation pays)
  \item \textbf{Extractive institutions:} $C_{ij} \lessgtr 0$ (rent-seeking, zero-sum)
\end{itemize}

\paragraph{Multiple Stable Equilibria.}
Both configurations can be stable:
\begin{center}
\renewcommand{\arraystretch}{1.2}
\begin{tabular}{lccc}
\hline
\textbf{Equilibrium} & $C_{ij}$ & $K$ & $Q$ \\
\hline
Inclusive & $> 0$ & $\approx 1$ & High \\
Extractive & $\lessgtr 0$ & $\approx 1$ & Low \\
\hline
\end{tabular}
\end{center}

\paragraph{Key Insight.}
High coherence ($K \approx 1$) does not imply high quality ($Q$).
A society can be stably trapped in a bad equilibrium.
This resolves the puzzle of persistent underdevelopment:
extractive equilibria are self-reinforcing.

\paragraph{Critical Junctures.}
Major shocks can shift $\Psi$ between equilibria.
The framework provides language for analyzing institutional transitions.

\subsection{Organizational Economics: The Supermodularity Foundation}

\paragraph{Roberts Organizational Design.}
\citet{roberts2004} models organizations as systems of complementary choices:
\begin{equation}
  \Pi = \Pi(\text{Strategy, Structure, Processes, People, Culture, Incentives})
\end{equation}
with $\partial^2 \Pi / \partial x_i \partial x_j > 0$ (supermodular complements).

\paragraph{Direct Correspondence.}
Roberts' complementarity matrix \emph{is} $C$:
\begin{equation}
  C^{Roberts}_{ij} = \frac{\partial^2 \Pi}{\partial x_i \partial x_j}
\end{equation}

Roberts' ``fit'' \emph{is} $K$:
\begin{equation}
  K = 1 \text{ at coherent configuration}, \quad K < 1 \text{ at misfit}
\end{equation}

\paragraph{Coherent Configurations.}
Due to complementarity, the profit function has multiple local maxima:
\begin{center}
\renewcommand{\arraystretch}{1.2}
\begin{tabular}{lllll}
\hline
\textbf{Configuration} & \textbf{Strategy} & \textbf{Structure} & \textbf{Culture} & \textbf{Incentives} \\
\hline
Cost Leader & Efficiency & Hierarchy & Discipline & Output-based \\
Innovator & Differentiation & Flat & Creative & Risk-reward \\
Customer Intimate & Service & Teams & Responsive & Satisfaction \\
\hline
\end{tabular}
\end{center}

Each is a peak ($K \approx 1$). Mixtures are valleys ($K < 1$).

\paragraph{Non-Concavity.}
Complementarity generates non-concave payoff landscapes.
This explains why partial organizational change often fails:
it lands in the valley between peaks.

\subsection{Financial Economics: The Market Coherence Extension}

\paragraph{CAPM.}
The Capital Asset Pricing Model relates expected returns to covariance:
\begin{equation}
  \mathbb{E}[r_i] = r_f + \beta_i (\mathbb{E}[r_m] - r_f)
\end{equation}

\paragraph{Covariance as Complementarity.}
In financial markets, the complementarity matrix is the return covariance matrix:
\begin{equation}
  C^{market}_{ij} = \text{Cov}(r_i, r_j)
\end{equation}

\paragraph{Market Coherence.}
\begin{equation}
  K_{market} = 1 - \frac{\|\Sigma_t - \Sigma^*\|_F}{\|\Sigma^*\|_F}
\end{equation}

\begin{itemize}
  \item $K_{market} \approx 1$: Correlations stable, diversification works, EMH approximately holds
  \item $K_{market} < 1$: Correlations break down, diversification fails, crisis regime
\end{itemize}

\paragraph{Financial Crises as Coherence Collapse.}
The 2008 crisis: $K_{market}$ dropped from $\approx 0.95$ to $\approx 0.3$.
``All assets fell together''---correlations spiked, destroying diversification.

\paragraph{Jensen's Alpha.}
Systematic $\alpha \neq 0$ indicates belief-payoff divergence:
agents' expectations don't match realized complementarities.
This is market incoherence.

\subsection{International Trade Theory: The Factor Endowment Extension}

\paragraph{Heckscher--Ohlin Model \citep{heckscher1919,ohlin1933}.}
The foundational model of international trade predicts that countries export goods
that use their abundant factor intensively. A capital-rich country exports
capital-intensive goods; a labor-rich country exports labor-intensive goods.

\paragraph{Standard Assumptions.}
\begin{enumerate}
  \item Identical technology across countries
  \item Identical preferences across countries
  \item Perfect factor mobility within, zero between countries
  \item No transport costs, perfect competition
\end{enumerate}

\paragraph{Translation to $(C, \Psi)$ Framework.}
Factor endowment is a component of context:
\begin{equation}
  \Psi_{factor} = (K/L) \quad \text{(Capital-Labor ratio)}
\end{equation}

Heckscher--Ohlin assumes all other context is identical:
\begin{align}
  \Psi_{tech}^A &= \Psi_{tech}^B \quad \text{(Same technology)} \\
  \Psi_{inst}^A &= \Psi_{inst}^B \quad \text{(Same institutions)} \\
  \Psi_{norm}^A &= \Psi_{norm}^B \quad \text{(Same preferences)}
\end{align}

Only $\Psi_{factor}^A \neq \Psi_{factor}^B$ varies.

\paragraph{Production Complementarity.}
The production function $Y = F(K, L)$ implies factor complementarity:
\begin{equation}
  C_{KL}^{prod} = \frac{\partial^2 F}{\partial K \partial L} > 0
\end{equation}
Capital and labor are complements in production.

\paragraph{Trade as Coherent Configuration.}
Given different factor endowments, there exists a coherent global specialization:
\begin{align}
  C^*(\Psi_{factor}^{high\ K/L}) &= \text{``Produce capital-intensive goods''} \\
  C^*(\Psi_{factor}^{low\ K/L}) &= \text{``Produce labor-intensive goods''}
\end{align}

Free trade equilibrium is the reference structure $C^*$ for the global system:
\begin{itemize}
  \item $K^{global} \approx 1$: Coherent international division of labor
  \item $Q^{trade} > Q^{autarky}$: Gains from trade
\end{itemize}

\paragraph{Heckscher--Ohlin as Limiting Case.}
\begin{equation}
  \text{H-O} = (C, \Psi)\text{-Framework} \Big|_{\substack{
    C_{ij}^{social} = 0 \\
    \Psi = \Psi_{factor} \text{ only} \\
    U = U^F \text{ only} \\
    \Psi_{t+1} = \Psi_t
  }}
\end{equation}

The restrictions are severe:
\begin{center}
\renewcommand{\arraystretch}{1.2}
\begin{tabular}{lll}
\hline
\textbf{H-O Assumption} & \textbf{Framework Translation} & \textbf{What Is Excluded} \\
\hline
Identical preferences & $U = U^F$ only & $U^S$, $U^{IDN}$ \\
Identical technology & $\Psi_{tech}$ constant & Innovation, tech transfer \\
Perfect competition & $C_{ij}^{market} = 0$ & Market power, networks \\
No institutions & $\Psi_{inst}$ irrelevant & Trade regimes, contracts \\
Static & $\Psi_{t+1} = \Psi_t$ & Path dependence, dynamics \\
\hline
\end{tabular}
\end{center}

\paragraph{What Heckscher--Ohlin Cannot Explain.}
The restrictions rule out:
\begin{itemize}
  \item \textbf{Protectionist politics:} If $U = U^F$ only, tariffs are irrational.
        But $U^{IDN}$ (``Made in America'') can outweigh $U^F$ losses.
  \item \textbf{Endogenous factor accumulation:} South Korea transformed from
        labor-abundant (1960) to capital/skill-abundant (2020).
        H-O explains the cross-section, not the transformation.
  \item \textbf{Institutional comparative advantage:} Hall--Soskice show that
        trade patterns depend on $\Psi_{inst}$, not just $\Psi_{factor}$.
  \item \textbf{Trade regime dynamics:} WTO erosion, bilateral deals, trade wars---
        these are $\Psi_{inst}^{global}$ dynamics that H-O ignores.
\end{itemize}

\paragraph{The Full Framework on Trade.}
Relaxing H-O restrictions yields richer predictions:

\textbf{1. Identity Utility in Trade Policy.}
\begin{equation}
  \Delta U^{total} = \Delta U^F + \omega_{IDN} \cdot \Delta U^{IDN}
\end{equation}
Tariffs may reduce $U^F$ but increase $U^{IDN}$ (``protecting our industries'').
For voters with high $\omega_{IDN}$, protectionism is rational.

\textbf{2. Endogenous Factor Accumulation.}
\begin{equation}
  \Psi_{factor,t+1} = f(\Psi_{factor,t}, \text{Investment}, \text{Education}, \text{Policy})
\end{equation}
Countries can \emph{change} their comparative advantage through policy.
Development strategy matters.

\textbf{3. Institutional Complementarity in Trade.}
\citet{hallsoskice2001} show:
\begin{itemize}
  \item LMEs (USA, UK): Comparative advantage in radical innovation
  \item CMEs (Germany, Japan): Comparative advantage in incremental innovation
\end{itemize}
This reflects $\Psi_{inst}$, not $\Psi_{factor}$.

\textbf{4. Context Hysteresis in Trade Wars.}
The Trump tariff analysis (Section~\ref{sec:predictions}) shows:
\begin{equation}
  \lim_{t \to \infty} \Psi_t \neq \Psi^{pre-tariff} \quad \text{even if tariffs removed}
\end{equation}
Tariffs shift $\Psi_{norm}$ (free trade norm weakens), $\Psi_{inst}$ (WTO erodes),
and $\Psi_{tech}$ (supply chains reconfigure). The path back is not the path forward.

\paragraph{Empirical Implications.}
The framework generates testable predictions beyond H-O:
\begin{enumerate}
  \item Trade policy support correlates with $U^{IDN}$ measures, not just $U^F$ exposure
  \item Countries with institutional complementarity trade more within type (LME-LME, CME-CME)
  \item Trade shocks have persistent effects even after policy reversal (hysteresis tests)
  \item Factor accumulation responds to institutional context, not just prices
\end{enumerate}

\paragraph{Policy Implications.}
\begin{itemize}
  \item Compensating losers with $U^F$ transfers may fail if $U^{IDN}$ losses dominate
  \item Trade agreements must address identity concerns, not just efficiency gains
  \item ``Adjustment assistance'' should target context change, not just retraining
  \item Global trade governance requires attention to $K^{global}$, not just $Q^{global}$
\end{itemize}

\subsection{Agency Theory: The Alignment Extension}

\paragraph{Jensen--Meckling.}
\citet{jensenmeckling1976} analyze principal-agent conflicts:
\begin{align}
  U_P &= f(e) - w && \text{(Principal wants effort)} \\
  U_A &= u(w) - c(e) && \text{(Agent dislikes effort)}
\end{align}

\paragraph{Agency Cost as Incoherence.}
Ideal alignment: $C_{PA}^* > 0$ (principal and agent interests complementary).
Actual alignment: $C_{PA}^{actual} < C_{PA}^*$ (misalignment).

\begin{equation}
  K_{PA} = 1 - \frac{\|C_{PA}^{actual} - C_{PA}^*\|}{\|C_{PA}^*\|}
\end{equation}

Agency costs measure the deviation from coherence.

\paragraph{Governance as Coherence Restoration.}
Corporate governance mechanisms (performance pay, monitoring, ownership concentration)
aim to increase $K_{PA}$ toward 1.

\subsection{Prospect Theory: The Loss Aversion Extension}

\paragraph{Kahneman--Tversky.}
\citet{kahnemantversky1979} showed that losses loom larger than gains:
\begin{equation}
  V = \sum_i \pi(p_i) \cdot v(x_i - r)
\end{equation}
with $|v(-x)| > v(x)$ for $x > 0$.

\paragraph{Integration in FEPSDE.}
The framework incorporates loss aversion directly:
\begin{equation}
  U = \sum_d \sum_t \delta_d^t [G_{d,t} - \lambda_d \cdot P_{d,t}]
\end{equation}
where $\lambda_d > 1$ captures loss aversion.

The separate treatment of Gains and Pains in the 144-component structure
embodies the Prospect Theory insight that positive and negative outcomes
are processed differently.

\subsection{Summary: The Theory Space}

Each theory occupies a position in ``theory space'' defined by restrictions:

\begin{center}
\renewcommand{\arraystretch}{1.2}
\small
\begin{tabular}{lccccc}
\hline
\textbf{Theory} & $C_{ij}$ & $\Psi$ dynamics & $U^{IDN}$ & $\lambda$ & Level \\
\hline
Arrow--Debreu & $= 0$ & Exogenous & No & $= 1$ & Individual \\
Heckscher--Ohlin & $= 0$ & Exogenous & No & $= 1$ & International \\
Fehr--Schmidt & $\neq 0$ & Slow & No & $= 1$ & Individual \\
Akerlof--Kranton & $\neq 0$ & Slow & \textbf{Yes} & $= 1$ & Individual \\
Bénabou--Tirole & $\neq 0$ & Slow & \textbf{Yes} & $= 1$ & Individual \\
Aghion--Howitt & $\neq 0$ & \textbf{Endogenous} & No & $= 1$ & Macro \\
Acemoglu--Robinson & $\neq 0$ & \textbf{Endogenous} & No & $= 1$ & Nation \\
Roberts & $> 0$ & Slow & No & $= 1$ & Organization \\
CAPM & $= \text{Cov}$ & Exogenous & No & $= 1$ & Market \\
Jensen--Meckling & $\lessgtr 0$ & Slow & No & $= 1$ & Organization \\
Prospect Theory & $\neq 0$ & Endogenous & No & $\mathbf{> 1}$ & Individual \\
\hline
\textbf{Full Framework} & Any & Endogenous & Yes & $> 1$ & Multi-level \\
\hline
\end{tabular}
\end{center}

\paragraph{The Unification Claim.}
The $(C, \Psi)$ framework is not merely another theory alongside others.
It is a \emph{meta-framework} that:
\begin{enumerate}
  \item Shows which restrictions generate which theory
  \item Explains why different theories apply in different contexts
  \item Reveals what each theory cannot explain (by identifying active restrictions)
  \item Provides a common language across economics subdisciplines
\end{enumerate}

Detailed derivations for each theory are in Appendix~\ref{app:limitingcases}.

%%%%%%%%%%%%%%%%%%%%%%%%%%%%%%%%%%%%%%%%%%%%%%%%%%%%%%%%%%%%%%%%%%%%%%%%%%%%%%%